% ===================================================================
% SKILL: SUNUM AKIŞI (TOC) VE BÖLÜM YAPISI
% ===================================================================
% TANIM:
% Bu dosya, sunumun ana omurgasını oluşturan iki kritik yapıyı içerir:
% 
% 1. OTOMATİK AKIŞ FRAME'İ: Her yeni bölüme geçildiğinde otomatik olarak
%    ekrana gelen ve sunumun neresinde olduğumuzu gösteren "Sunum Akışı"
%    ekranının kodudur. İçeriği çok olan derslerde "Çift Sütun" (Double Column)
%    yapısı kullanılarak taşma engellenir.
% 
% 2. STANDART İÇERİK İSKELETİ: İçeriğinizi oluştururken kullanmanız gereken
%    Section > Subsection > Frame hiyerarşisinin boş şablonudur.
% ===================================================================

% -------------------------------------------------------------------
% 1. OTOMATİK SUNUM AKIŞI (PREAMBLE KISMI)
% -------------------------------------------------------------------
% Bu kod bloğunu \begin{document} öncesine yapıştırın.
% Her \section{} komutu algılandığında bu frame otomatik üretilir.

\AtBeginSection[]{
  \begin{frame}{Sunum Akışı} % Frame Başlığı
  \footnotesize % Yazıları sığdırmak için küçültülmüş font
  
  % Maddeler arası boşluğu ayarlar (daha şık görünüm için)
  \addtobeamertemplate{section in toc}{}{\vspace{0.15cm}}
  
  % Çift Sütun Yapısı (Eğer bölüm sayısı az ise tek sütun yapabilirsiniz)
  \begin{columns}[t]
    % SOL SÜTUN (İlk 3 Bölüm)
    \column{.5\textwidth}
    \tableofcontents[sections={1--3},currentsection]
    
    % SAĞ SÜTUN (Kalan Bölümler)
    \column{.5\textwidth}
    \tableofcontents[sections={4--10},currentsection]
  \end{columns}
  \end{frame}
}

% -------------------------------------------------------------------
% 2. İÇERİK İSKELETİ (DOCUMENT BODY KISMI)
% -------------------------------------------------------------------
% Sunumunuzu oluştururken aşağıdaki hiyerarşiyi kopyalayıp çoğaltın.

% ANA BÖLÜM (Otomatik Akış Ekranını Bu Tetikler)
\section{Bölüm Başlığı} 

    % ALT BÖLÜM (İsteğe bağlı, karmaşık konuları bölmek için)
    \subsection{Alt Konu Başlığı}

        % STANDART BOŞ FRAME
        \begin{frame}{Slayt Başlığı}
            \justifying{} % Metni iki yana yaslar (Okunabilirlik için önemli)
            
            % İçerik buraya gelecek...
            Buraya metin, madde işaretleri veya görseller gelecek.
            
        \end{frame}

        % İKİNCİ FRAME ÖRNEĞİ
        \begin{frame}{Slayt Başlığı 2}
            \begin{itemize}
                % İçerik buraya gelecek...
                Buraya metin, madde işaretleri veya görseller gelecek.
            \end{itemize}
        \end{frame}
